\section*{Chapitre \ref{ch:b2:elements-in-the-earth} --- Les éléments de la Terre}

\begin{solution}{exo:b2:elements-in-the-earth:1}
Le potassium représente $2\times39.10/(2\times39.10+16.00) = 78.20/94.20 = 0.8301$ de la
masse de \ce{K2O}, donc $\num{5130}\times0.8301 = \qty{4259}{ppm}$ de K. Le titane représente
$47.87/79.87 = 0.5993$ de \ce{TiO2}~: $\num{27310}\times0.5993 = \qty{16370}{ppm}$ de Ti.
\end{solution}

\begin{solution}{exo:b2:elements-in-the-earth:2}
L'aluminium est lithophile et réfractaire~: la Terre silicatée en a tout conservé, dans des
proportions proches de celles des \omterm{def:b2:planet-formation-meteorites:chondrule}{chondrites} (1.16). Le sodium est lithophile mais volatil
(\qty{958}{K})~: il a été perdu avant ou pendant l'accrétion, ne laissant que 0.22. Le
nickel se condense avec le magnésium (\qty{1353}{K})~; la volatilité ne peut donc pas
expliquer 0.077~: il est sidérophile, et l'essentiel est passé dans le noyau. L'iridium est
réfractaire (\qty{1603}{K}) et fortement sidérophile~: 0.0029, presque tout dans le noyau.
\end{solution}

\begin{solution}{exo:b2:elements-in-the-earth:3}
$\mathrm{La}_N = 31/0.232 = 133.6$, $\mathrm{Yb}_N = 2.0/0.163 = 12.27$, et
$(\mathrm{La}/\mathrm{Yb})_N = 10.9$. La croûte est enrichie en toutes les \omterm{def:b2:elements-in-the-earth:ree}{terres
rares}, et onze fois plus en les légères qu'en les lourdes~: elle est faite de
liquides, qui prennent le plus volontiers les \omterm{def:b2:elements-in-the-earth:ree}{terres rares} légères incompatibles.
\end{solution}

\begin{solution}{exo:b2:elements-in-the-earth:4}
Pour le samarium, $\mathrm{Sm}_N = 4.7/0.145 = 32.4$, $\mathrm{Eu}_N = 1.0/0.0546 = 18.3$,
$\mathrm{Gd}_N = 4.0/0.198 = 20.2$~; $\mathrm{Eu}^* = \sqrt{32.4\times20.2} = 25.6$
et $\mathrm{Eu}/\mathrm{Eu}^* = 18.3/25.6 = 0.72$, une anomalie négative.
Le plagioclase, qui intègre \ce{Eu^2+}, a cristallisé à partir des liquides
qui ont construit la croûte supérieure et qui sont restés plus profonds.
\end{solution}

\begin{solution}{exo:b2:elements-in-the-earth:5}
$\bar D_{\mathrm{La}} = 0.65\times0.058+0.18\times0.001+0.17\times0.001 = 0.0377 +
0.0002 + 0.0002 = 0.038$~; $\bar D_{\mathrm{Lu}} = 0.3445+1.044+0.0306 = 1.42$. Le lutécium est
compatible ($\bar D > 1$), le lanthane fortement incompatible. Le grenat porte
$1.044/1.42 = 74$ pour cent du lutécium de la roche,
alors qu'il ne représente que 18 pour cent de la roche.
\end{solution}

\begin{solution}{exo:b2:elements-in-the-earth:6}
Pour le sodium, le rapport Na/Mg de la Terre silicatée vaut $2670/\num{228000} = 0.01171$ et la valeur CI $5010/\num{95870} = 0.05226$~: le rapport vaut 0.224. Pour le nickel, $1960/\num{228000}
= 0.008596$ contre $\num{10640}/\num{95870} = 0.1110$~: 0.0775. Le sodium est
appauvri par sa volatilité (\qty{958}{K})~; le nickel se condense
au-dessus de \qty{1300}{K} et il est trois fois plus appauvri que le sodium
parce qu'il est sidérophile~:
le noyau en a pris la plus grande part.
\end{solution}

\begin{solution}{exo:b2:elements-in-the-earth:7}
Pour le lanthane~: $5.21\times0.058 = \qty{0.302}{ppm}$~; pour le lutécium~: $0.53\times0.53 = \qty{0.281}{ppm}$.
Les valeurs de normalisation s'annulent dans le rapport des rapports~:
$(\mathrm{La}/\mathrm{Lu})_{N,\mathrm{cpx}}/(\mathrm{La}/\mathrm{Lu})_{N,\mathrm{melt}} =
D_{\mathrm{La}}/D_{\mathrm{Lu}} = 0.058/0.53 = 0.11$~: le motif du cristal est neuf fois plus pentu
vers les \omterm{def:b2:elements-in-the-earth:ree}{terres rares} lourdes que celui du liquide.
\end{solution}

\begin{solution}{exo:b2:elements-in-the-earth:8}
Pour les trois réservoirs, le rapport normalisé Nb/La vaut successivement $(12/0.658)/(31/0.648) = 18.24/47.84 = 0.38$ pour la croûte,
$(18.10/0.658)/(15.20/0.648) = 27.51/23.46 = 1.17$ pour le basalte hawaïen et $7.96/8.04 = 0.99$ pour le MORB.
Seule la croûte, au fort déficit en niobium, porte la signature de la subduction.
\end{solution}

\begin{solution}{exo:b2:elements-in-the-earth:9}
\omterm{def:b2:elements-in-the-earth:bse}{Manteau primitif}~: $240/0.0203 = \num{11800}$. MORB~: K $= 0.160\times0.830\times10^4 =
\qty{1328}{ppm}$, K/U $= \num{11200}$. Croûte~: K $= 2.80\times0.830\times10^4 =
\qty{23240}{ppm}$, K/U $= \num{8600}$. Le potassium est
multiplié par cent environ entre le manteau et la croûte, pourtant K/U ne change que de moins d'un tiers~: les deux éléments sont à peu près aussi incompatibles,
et la fusion les sépare à peine. L'uranium est réfractaire,
si bien que son abondance découle du modèle chondritique~; K/U mesuré sur des laves
donne alors le potassium de la planète, qui, étant volatil, ne peut pas être pris dans les \omterm{def:b2:planet-formation-meteorites:chondrule}{chondrites}.
\end{solution}

\begin{solution}{exo:b2:elements-in-the-earth:10}
À l'équilibre, $C_{\mathrm{rock}} = \bar D\,C_{\mathrm{melt}}$ pour chaque élément (et une petite fraction de liquide
laisse la composition de la roche presque inchangée), donc $(C_1/C_2)_{\mathrm{melt}} = (C_{1}/\bar D_1)/(C_{2}/\bar D_2)$ pour les $C_1,
C_2$ de la roche, et le rapport vaut $\bar D_2/\bar D_1$.
Pour le lanthane et le lutécium~: $1.419/0.03805 = 37$.
Le liquide d'un manteau à grenat est 37 fois plus riche en lanthane par rapport au lutécium que sa source~:
le grenat a retenu les \omterm{def:b2:elements-in-the-earth:ree}{terres rares} lourdes, et le motif du liquide est abrupt.
\end{solution}

\begin{solution}{exo:b2:elements-in-the-earth:11}
Pour chaque élément, $C_{\mathrm{mix}} = fC_U + (1-f)C_M$~: linéaire en $f$.
Le rapport $\bigl(f\,\mathrm{Th}_U + (1-f)\mathrm{Th}_M\bigr)/\bigl(f\,\mathrm{Nb}_U +
(1-f)\mathrm{Nb}_M\bigr)$ est un quotient de deux fonctions linéaires,
une hyperbole en $f$. Th/Nb $= 0.2$ lorsque $f(10.5-2.4) + (1-f)(0.404-1.048) = 0$, soit $8.1f = 0.644(1-f)$,
$f = 0.644/8.744 = 0.074$~: environ 7 pour cent de croûte font passer Th/Nb de 0.077 à 0.2.
Pour $f = 0.1$~: Th $= 1.05 + 0.364 = 1.414$, Nb $= 1.2 + 4.716 = 5.916$, Th/Nb $= 0.24$.
\end{solution}

\begin{solution}{exo:b2:elements-in-the-earth:12}
Retirer un composant qui ne contient ni l'un ni l'autre élément laisse les deux masses, donc
leur rapport, inchangé. Si tout l'aluminium a été conservé, $(\mathrm{Al}/\mathrm{Mg}) = 1.16$
fois la valeur CI signifie que la Terre silicatée contient $1/1.16 = 0.86$ de son
magnésium chondritique~: 14 pour cent manquent. Le silicium rapporté à l'aluminium vaut
$0.83/1.16 = 0.72$ fois la valeur CI~: 28 pour cent du silicium chondritique manqueraient.
\end{solution}

\begin{solution}{pb:b2:elements-in-the-earth:1}
\textbf{1.} Fe/Mg $= \num{182800}/\num{95870} = 1.907$~; Ni/Mg $= 0.1110$~; Si/Mg $=
1.111$~; Al/Mg $= 0.0887$.

\textbf{2.} $(\num{182800}/55.85)/(\num{106500}/28.09) = 3273/3791 = 0.863$ atome de fer
par atome de silicium, à 7 pour cent près du 0.81 de la photosphère~: raisonnable,
puisque les \omterm{def:b2:elements-in-the-earth:ci}{chondrites CI} reproduisent le Soleil pour les éléments formateurs de roches.

\textbf{3.} La perte d'un composant qui ne porte aucun des deux éléments d'un rapport laisse ce rapport
inchangé (\cref{exo:b2:elements-in-the-earth:12})~: seules les concentrations absolues changent.

\textbf{4.} Le magnésium est lithophile et n'entre pas dans le noyau~; c'est un élément
majeur, bien déterminé à la fois dans les \omterm{def:b2:planet-formation-meteorites:chondrule}{chondrites} et dans le manteau. Le fer est l'élément
qui se partage entre noyau et manteau~: il ne peut pas servir de référence à lui-même.

\textbf{5.} Des éléments qui se condensent ensemble n'ont pas été séparés par leur volatilité dans la
nébuleuse, si bien que la Terre globale devrait conserver leurs rapports. Al/Mg s'en écarte~: l'aluminium
est réfractaire (\qty{1653}{K}), et le rapport Al/Mg de la Terre silicatée vaut 1.16 fois la valeur CI.

\textbf{6.} Fe $= 8.05\times55.85/71.85 = 6.26$ pour cent~; Ni $= 0.25\times58.69/74.69 =
0.196$ pour cent.

\textbf{7.} Fe/Mg $= 6.26/22.8 = 0.2746$~; Ni/Mg $= 0.196/22.8 = 0.00860$.

\textbf{8.} Fer~: $1 - 0.2746/1.907 = 0.856$, soit 85.6 pour cent
manquant~; nickel~: $1 - 0.00860/0.1110 = 0.923$, soit 92.3 pour cent.

\textbf{9.} L'iridium est plus fortement sidérophile que le nickel, et le nickel plus que
le fer~: plus un élément préfère le métal, plus le métal en a entraîné vers le noyau.

\textbf{10.} Dans les silicates, sous forme de \ce{Fe^2+} dans l'olivine et les
pyroxènes (8 pour cent de FeO)~: le fer oxydé en FeO se lie à l'oxygène comme un
\omterm{def:b2:elements-in-the-earth:goldschmidt}{élément lithophile} et est resté dans le manteau~; seul le fer métallique a coulé.

\textbf{11.} $(1-x)C_{\mathrm{Fe}} + xc_{\mathrm{Fe}} = r_{\mathrm{Fe}}(1-x)C_{\mathrm{Mg}}$~;
en divisant par $(1-x)c_{\mathrm{Fe}}$, on obtient la proposition.

\textbf{12.} $r_{\mathrm{Fe}}C_{\mathrm{Mg}} - C_{\mathrm{Fe}} = 1.907\times22.8 - 6.26 =
37.22$~; $y = 37.22/85.5 = 0.4353$, $x = y/(1+y) = 0.303$.

\textbf{13.} $0.1110\times22.8 - 0.196 = 2.335$~; $y = 2.335/5.2 = 0.4490$, $x = 0.310$.

\textbf{14.} $y = (85.5\times37.22 + 5.2\times2.335)/(85.5^2 + 5.2^2) =
3194.4/7337.3 = 0.4354$, $x = 0.303$.

\textbf{15.} Chaque équation pèse comme le carré de son coefficient~:
$85.5^2/5.2^2 = 270$ fois plus pour le fer, l'élément de beaucoup le plus abondant.

\textbf{16.} $37.22/2.335 = 15.9$, contre
$85.5/5.2 = 16.4$ dans le modèle~: à 3 pour cent près.

\textbf{17.} $0.303\times\num{5.972e24} = \qty{1.81e24}{kg}$.

\textbf{18.} $(32.5-30.3)/32.5 = 7$ pour cent plus faible.

\textbf{19.} $y = 0.325/0.675 = 0.4815$, Fe dans le noyau $= 37.22/0.4815 = 77.3$ pour cent,
et Ni $= 2.335/0.4815 = 4.85$ pour cent~: 18 pour cent d'éléments légers, soit environ le double
des 5 à 10 pour cent permis par la masse volumique. Le rapport Fe/Mg de la Terre doit être
un peu plus élevé que la valeur CI, ou son magnésium silicaté plus élevé que celui du modèle.

\textbf{20.} Silicium, soufre, carbone, oxygène et hydrogène. Le silicium manque dans la Terre
silicatée~: son rapport normalisé au magnésium, 0.83, est inférieur au niveau réfractaire.

\textbf{21.} $1.111\times22.8 - 21.0 = 4.331$~;
$4.331/0.4354 = 9.9$ pour cent de Si dans le noyau.

\textbf{22.} 9.9 contre 6 pour cent~: le noyau pourrait contenir plus de silicium que le modèle,
ou une partie du silicium n'a jamais été accrétée. Le silicium se condense à \qty{1310}{K}, un peu
en dessous du magnésium et du fer (1336 et \qty{1334}{K})~: une partie a pu rester dans le gaz.

\textbf{23.} Noyau~: $0.325\times0.5 = \qty{0.1625}{ppm}$ de la masse
de la Terre~; Terre silicatée~: $0.675\times0.001 = \qty{0.000675}{ppm}$.
Le noyau contient $0.1625/0.1632 = 99.6$ pour cent de l'or.

\textbf{24.} La chimie seule place
\textbf{\qty[text-series-to-math=true]{30.3}{\percent}} de la masse de la Terre dans son noyau,
à 7 pour cent près des 32.5 pour cent sismologiques.
\end{solution}
